\labsection{11}{Hashing and collision resolution}{sec:lab11}

\begin{sourcealignment}{Official Lab Manual --- Lab XI}
\textbf{Objective.} Learn hashing techniques for efficient storage and access.

\textbf{Outcome.} Store data so it can be accessed efficiently.

\textbf{Problem directions.} Practice hashing and resolve collisions using linear probing, double hashing, and quadratic hashing/probing.
\end{sourcealignment}

\begin{labproject}{Resolve the same collision three ways}
\textbf{Coverage.} Chapter 11.\\
\textbf{Source files.} \texttt{Lab09/Task01.cpp} --- linear probing, quadratic probing, and double hashing on one table.\\
\textbf{Goal.} Resolve the same collisions three ways --- linear probing, quadratic probing, and double hashing --- and record every probe.

Use the source example $H(x)=x\bmod5$ with keys 10, 12, and 14, then insert 15 so the home slot 0 collides with 10.

\begin{lstlisting}
vector<int> keys = {10, 12, 14, 15};
const int TABLE_SIZE = 5;
\end{lstlisting}

\begin{tryit}
For each collision strategy, record every probe index rather than only the final slot. Then add more keys and compare how probe sequences grow as the table becomes crowded.
\end{tryit}
\end{labproject}

\begin{verification}
\begin{itemize}
  \item Every inserted key remains discoverable using the same probe rule used for insertion.
  \item Probe sequences stay inside the table using modular arithmetic.
  \item The double-hash step does not collapse to zero and is suitable for the chosen table size.
  \item The term ``quadratic hashing'' in the source is connected to the usual term ``quadratic probing.''
\end{itemize}
\end{verification}

\begin{labdeliverable}
Submit all three collision-resolution methods and a comparison table listing the probe sequence and final slot for each key.
\end{labdeliverable}
